\documentclass[11pt,a4paper]{article}

\usepackage[utf8]{inputenc}
\usepackage[T1]{fontenc}
\usepackage{geometry}
\geometry{margin=2.5cm}
\usepackage{listings}
\usepackage{xcolor}
\usepackage{hyperref}
\usepackage{enumitem}
\usepackage{booktabs}
\usepackage{parskip}
\usepackage{titlesec}

\hypersetup{colorlinks=true, linkcolor=blue!50!black, urlcolor=blue!50!black, citecolor=blue!50!black}

\definecolor{codebg}{RGB}{245,245,245}
\definecolor{codegray}{RGB}{110,110,110}
\definecolor{codegreen}{RGB}{0,110,0}
\definecolor{codeblue}{RGB}{30,80,180}

\lstdefinestyle{code}{
    backgroundcolor=\color{codebg}, basicstyle=\ttfamily\footnotesize,
    keywordstyle=\color{codeblue}, commentstyle=\color{codegray}\itshape,
    stringstyle=\color{codegreen}, breaklines=true, breakatwhitespace=false,
    frame=single, rulecolor=\color{black!15}, numbers=left,
    numberstyle=\tiny\color{codegray}, xleftmargin=1.8em, framexleftmargin=1.5em,
    showstringspaces=false, tabsize=2,
    literate=
      {á}{{\'a}}1 {à}{{\`a}}1 {ã}{{\~a}}1 {â}{{\^a}}1
      {é}{{\'e}}1 {ê}{{\^e}}1 {í}{{\'i}}1
      {ó}{{\'o}}1 {õ}{{\~o}}1 {ô}{{\^o}}1 {ú}{{\'u}}1
      {ç}{{\c c}}1 {Á}{{\'A}}1 {É}{{\'E}}1 {Í}{{\'I}}1
      {Ó}{{\'O}}1 {Ú}{{\'U}}1 {Ã}{{\~A}}1 {Õ}{{\~O}}1
      {Ç}{{\c C}}1,
}
\lstset{style=code}

\titleformat{\section}{\normalfont\Large\bfseries}{\thesection}{1em}{}
\titleformat{\subsection}{\normalfont\large\bfseries}{\thesubsection}{1em}{}

\sloppy

\title{\textbf{TelemetriaMqtt v2 --- Checkpoint 6} \\[4pt]
\large Dashboards Grafana --- caminho histórico}
\author{Projeto de Telemetria FSAE}
\date{16 de setembro de 2026}

\begin{document}
\maketitle

\begin{abstract}
\noindent Este documento detalha o sexto checkpoint do projeto: colocar o Grafana no ar como container Docker, com o datasource InfluxDB e um dashboard de análise histórica provisionados inteiramente como código -- nenhuma configuração manual pela interface web.
\end{abstract}

\tableofcontents

\section{Contexto e objetivo}

O caminho histórico da arquitetura (definida no \texttt{SCOPE.md}) termina em dashboards Grafana consultando o InfluxDB (colocado no ar no checkpoint 5). A decisão de arquitetura relevante aqui: assim como o broker e o banco, o Grafana também é provisionado \textbf{como código} -- datasource e dashboard são arquivos versionados no repositório, não cliques numa interface que ninguém mais consegue reproduzir. É a mesma razão pela qual o projeto rejeitou Node-RED no planejamento inicial.

\section{Estrutura adicionada}

\begin{lstlisting}[language=bash,caption={Arquivos novos}]
infra/
  docker-compose.yml                      (+ serviço grafana)
  grafana/
    generate_dashboard.py                  (gera o dashboard a partir de uma lista)
    dashboards/
      historico.json                       (gerado, não editado à mão)
    provisioning/
      datasources/influxdb.yml
      dashboards/dashboards.yml
\end{lstlisting}

\subsection{Datasource}

\begin{lstlisting}[language=yaml,caption={infra/grafana/provisioning/datasources/influxdb.yml}]
datasources:
  - name: InfluxDB
    uid: influxdb-main
    type: influxdb
    access: proxy
    url: http://influxdb:8086
    isDefault: true
    jsonData:
      version: Flux
      organization: $INFLUXDB_INIT_ORG
      defaultBucket: $INFLUXDB_INIT_BUCKET
    secureJsonData:
      token: $INFLUXDB_INIT_ADMIN_TOKEN
\end{lstlisting}

Dois pontos relevantes: a linguagem de consulta é \textbf{Flux} (não InfluxQL, a linguagem legada), consistente com o InfluxDB 2.x; e as credenciais vêm das \textbf{mesmas} variáveis de ambiente já usadas pelo próprio serviço InfluxDB (\texttt{infra/.env}) -- o Grafana suporta expansão de variável de ambiente diretamente nos arquivos de provisionamento, então não existe uma segunda cópia das credenciais em lugar nenhum. O \texttt{uid} fixo (\texttt{influxdb-main}) permite que o dashboard referencie esse datasource diretamente, sem precisar descobrir o UID em tempo de execução.

\subsection{Dashboard gerado por script}

Em vez de editar o JSON do dashboard à mão (o schema do Grafana é grande e repetitivo -- cada painel tem dezenas de campos), um script Python (\texttt{generate\_dashboard.py}) gera o arquivo a partir de uma lista declarativa de painéis:

\begin{lstlisting}[language=Python,caption={Trecho de generate\_dashboard.py}]
PANELS = [
    {"title": "RPM", "fields": ["rpm"], "unit": "none"},
    {"title": "Acelerador / Lambda / MAP",
     "fields": ["throttlePosition", "lambda", "map"], "unit": "none"},
    {"title": "Temperaturas (°C)",
     "fields": ["engineTemp", "airTemp"], "unit": "celsius"},
    {"title": "Correntes (A)",
     "fields": ["ventoinhaCurrent", "bombaCurrent", ...], "unit": "amp"},
    {"title": "Velocidade GPS (km/h)", "fields": ["gpsSpeed"], ...},
    {"title": "Tensão da bateria (V)", "fields": ["extBattery"], ...},
    {"title": "Estados digitais",
     "fields": ["bombaInput", "giratoriaInput"], "unit": "none"},
]
\end{lstlisting}

Os 20 sinais de telemetria (mesma lista do \texttt{mock\_publisher}/firmware) foram agrupados em 7 painéis por categoria, em vez de 20 painéis individuais -- mais legível pra uma primeira análise histórica. Cada painel vira uma consulta Flux gerada automaticamente:

\begin{lstlisting}[caption={Query Flux gerada pro painel "RPM"}]
from(bucket: "telemetria")
  |> range(start: v.timeRangeStart, stop: v.timeRangeStop)
  |> filter(fn: (r) => r._measurement == "telemetry")
  |> filter(fn: (r) => r._field == "rpm")
\end{lstlisting}

\textbf{Acoplamento a documentar}: o nome do bucket (\texttt{"telemetria"}) está fixo nas queries geradas -- o Flux não herda automaticamente o \texttt{defaultBucket} configurado no datasource. Se o bucket for renomeado no futuro, o dashboard precisa ser regenerado (\texttt{python3 generate\_dashboard.py > dashboards/historico.json}), não só o datasource reconfigurado.

\section{Testes}

Quatro testes validam este checkpoint contra um Grafana real (não mockado):

\begin{enumerate}[label=\textbf{Teste \arabic*}]
    \item \texttt{test\_docker\_compose\_up\_starts\_a\_working\_grafana} -- sobe influxdb+grafana via \texttt{docker-compose.yml} real, confirma: o healthcheck responde; o datasource provisionado tem o tipo, organização e bucket corretos; o dashboard \texttt{telemetria-historico} aparece na busca com exatamente 7 painéis.
    \item \texttt{test\_grafana\_queries\_real\_data\_through\_provisioned\_datasource} -- a validação mais forte do checkpoint: escreve um ponto no InfluxDB usando o \texttt{make\_influx\_writer} de produção, e consulta de volta através da \textbf{própria API do Grafana} (\texttt{/api/ds/query}, o mesmo endpoint que os painéis do dashboard usam internamente) -- não direto no InfluxDB. Confirma que a rede entre os containers do compose, a autenticação do datasource e a linguagem Flux configurada funcionam de ponta a ponta, não só que "a configuração parece certa" olhando os arquivos.
\end{enumerate}

\subsection{Resultado da execução}

\begin{lstlisting}[caption={Saída real de pytest -v}]
test_docker_compose.py::test_docker_compose_up_starts_a_working_grafana PASSED
test_docker_compose.py::test_grafana_queries_real_data_through_provisioned_datasource PASSED
\end{lstlisting}

\section{Validação manual}

Antes mesmo de escrever os testes automatizados, o comportamento foi confirmado manualmente via \texttt{curl} contra a API REST do Grafana: \texttt{/api/health}, \texttt{/api/datasources} (confirmando organização/bucket corretos), \texttt{/api/search} e \texttt{/api/dashboards/uid/telemetria-historico} (confirmando os 7 painéis), e por fim uma consulta real via \texttt{/api/ds/query} depois de publicar dados reais através do pipeline completo (mock publisher $\to$ broker $\to$ ingestão), retornando os pontos de RPM esperados.

\section{Nota sobre metodologia}

Este checkpoint é o exemplo mais claro, até agora, de uma inversão da ordem de TDD no projeto: toda a infraestrutura (serviço no compose, datasource, dashboard) foi construída primeiro, validada manualmente via \texttt{curl}, e só depois os testes automatizados foram escritos -- essencialmente codificando um comportamento que já tinha sido confirmado, em vez de guiar a implementação. Identificado numa revisão de processo com o usuário logo após este checkpoint ser fechado. A decisão tomada foi não refazer os checkpoints 5 e 6 (a cobertura de teste já é real, contra infraestrutura de verdade, e passa) -- mas a partir do checkpoint 7 o teste volta a vir primeiro, inclusive para infraestrutura declarativa como esta.

\section{Resumo do que ficou pronto}

\begin{itemize}
    \item Grafana (\texttt{grafana-oss:13.0.2}) real, provisionado inteiramente como código.
    \item Datasource InfluxDB (Flux) com credenciais compartilhadas via variável de ambiente, sem duplicação de segredo.
    \item Dashboard de 7 painéis, gerado por script a partir de uma lista declarativa, cobrindo os 20 sinais de telemetria.
    \item Validação de ponta a ponta confirmada: dado escrito no InfluxDB é visível através da própria API de consulta do Grafana.
\end{itemize}

\section{Próximos passos}

Checkpoint 7: dashboard Grafana ao vivo, usando Grafana Live + o plugin de datasource MQTT, pro caminho de visualização em tempo real (sem depender do InfluxDB) definido na arquitetura original do projeto.

\end{document}
